% Invoice - Pressa public template (MIT). https://pressa.dev/templates/invoice
% Line totals, subtotal, tax and total are computed from the data, each rounded to the
% currency's minor unit (cents, whole yen) as printed, so the printed numbers add up.
% Unit prices keep up to 4 decimals for items priced below a cent.
\documentclass[10pt,a4paper]{article}
\usepackage[T1]{fontenc}
\usepackage[default]{sourcesanspro}
\usepackage[margin=18mm,top=14mm]{geometry}
\usepackage[table]{xcolor}
\usepackage{tabularx}
\usepackage{longtable}
\usepackage{array}
\usepackage[hidelinks]{hyperref}

\definecolor{brand}{HTML}{4338CA}
\definecolor{ink}{HTML}{1F2937}
\definecolor{muted}{HTML}{6B7280}
\definecolor{band}{HTML}{EEF2FF}
\pagestyle{empty}
\setlength{\parindent}{0pt}
\renewcommand{\arraystretch}{1.55}
\newcolumntype{R}{>{\raggedleft\arraybackslash}X}
\newcolumntype{L}{>{\raggedright\arraybackslash}X}
\color{ink}

\begin{document}

{% assign d = currency | currency_decimals %}{% assign subtotal = 0 %}{% for item in items %}{% assign price = item.unit_price | round: 4 %}{% assign line = item.quantity | times: price | round: d %}{% assign subtotal = subtotal | plus: line %}{% endfor %}{% assign tax = subtotal | times: tax_rate | divided_by: 100.0 | round: d %}{% assign total = subtotal | plus: tax %}
\noindent\colorbox{brand}{\parbox{\dimexpr\linewidth-2\fboxsep}{\vspace{9pt}
  \color{white}
  \begin{tabularx}{\linewidth}{@{\hspace{6pt}}L R@{\hspace{6pt}}}
    {\LARGE\bfseries {{ seller.name }}} & {\fontsize{26}{28}\selectfont\bfseries INVOICE} \\
    {{ seller.address }} & No. {{ invoice_number }} \\
    {{ seller.email }} & {{ issue_date }}
  \end{tabularx}\vspace{9pt}}}

\vspace{16pt}
{\renewcommand{\arraystretch}{1.25}
\begin{tabularx}{\linewidth}{@{}L L R@{}}
  {\small\color{muted}BILL TO} & {\small\color{muted}PAYMENT DUE} & {\small\color{muted}AMOUNT DUE} \\[-4pt]
  {\bfseries {{ client.name }}} & {\bfseries {{ due_date }}} & {\Large\bfseries\color{brand} {{ total | currency: currency }}} \\[-4pt]
  {{ client.address }} & {{ payment_terms }} & {\small\color{muted}{{ currency }}}
\end{tabularx}}

\vspace{18pt}
% A longtable, so a long list breaks across pages with the header repeated.
% Every column has a fixed width: one compile pass gives aligned columns.
{\rowcolors{2}{band}{white}\setlength{\LTpre}{0pt}\setlength{\LTpost}{0pt}\setlength{\LTleft}{0pt}\setlength{\LTright}{0pt}
\begin{longtable}{@{\hspace{6pt}}>{\raggedright\arraybackslash}p{\dimexpr\linewidth-12pt-6\tabcolsep-12mm-24mm-24mm\relax} >{\raggedleft\arraybackslash}p{12mm} >{\raggedleft\arraybackslash}p{24mm} >{\raggedleft\arraybackslash}p{24mm}@{\hspace{6pt}}}
  \rowcolor{ink}\textcolor{white}{\textbf{Description}} & \textcolor{white}{\textbf{Qty}} & \textcolor{white}{\textbf{Unit price}} & \textcolor{white}{\textbf{Amount}} \\
\endhead
{% for item in items %}{% assign price = item.unit_price | round: 4 %}{% assign line = item.quantity | times: price | round: d %}  {{ item.description }} & {{ item.quantity | round: 2 }} & {{ price | currency: currency, 4 }} & {{ line | currency: currency }} \\
{% endfor %}\end{longtable}}

\vspace{6pt}
\hfill\begin{tabular}{@{}l r@{\hspace{6pt}}}
  {\color{muted}Subtotal} & {{ subtotal | currency: currency }} \\
  {\color{muted}Tax ({{ tax_rate }}\%)} & {{ tax | currency: currency }} \\
  \arrayrulecolor{brand}\hline
  \textbf{Total} & \textbf{\color{brand} {{ total | currency: currency }}} \\
\end{tabular}

\vfill
{\color{brand}\rule{\linewidth}{1pt}}\\[4pt]
\begin{tabularx}{\linewidth}{@{}L L@{}}
  {\small\color{muted}PAYMENT DETAILS} & {\small\color{muted}NOTES} \\[-3pt]
  {\small {{ payment_details }}} & {\small {{ notes }}}
\end{tabularx}

\end{document}
